% Versão breve (5-8 páginas) – compilar com XeLaTeX (xelatex artigo_breve.tex, duas vezes)
\documentclass[12pt,a4paper]{article}
\usepackage{fontspec}
\setmainfont{TeX Gyre Termes}
\usepackage{polyglossia}
\makeatletter
\ifcsname l@portuges\endcsname\else\expandafter\let\csname l@portuges\endcsname\l@english\fi
\ifcsname l@portuguese\endcsname\else\expandafter\let\csname l@portuguese\endcsname\l@english\fi
\makeatother
\setdefaultlanguage[variant=brazilian]{portuguese}
\usepackage[top=3cm,left=3cm,bottom=2cm,right=2cm]{geometry}
\usepackage{setspace}
\onehalfspacing
\usepackage{indentfirst}
\setlength{\parindent}{1.25cm}
\setlength{\parskip}{0pt}
\emergencystretch=1.5em
\usepackage{amsmath}
\usepackage{graphicx}
\graphicspath{{figuras/}}
\usepackage{float}
\usepackage{booktabs}
\usepackage{array}
\usepackage{caption}
\captionsetup{font=small,labelsep=endash,justification=centering,skip=4pt}
\usepackage{titlesec}
\titleformat{\section}{\normalsize\bfseries}{\thesection}{1em}{\MakeUppercase}
\titleformat{\subsection}{\normalsize}{\thesubsection}{1em}{\MakeUppercase}
\titlespacing*{\section}{0pt}{16pt}{10pt}
\titlespacing*{\subsection}{0pt}{10pt}{10pt}
\usepackage{fancyhdr}
\pagestyle{fancy}
\fancyhf{}
\fancyhead[R]{\small\thepage}
\renewcommand{\headrulewidth}{0pt}
\usepackage[hidelinks]{hyperref}
\urlstyle{same}
\newcommand{\fonte}[1]{\par\vspace{2pt}{\centering\small #1\par}\vspace{6pt}}
\newenvironment{referencias}{\begin{singlespace}\setlength{\parindent}{0pt}\setlength{\parskip}{6pt}\raggedright\footnotesize}{\end{singlespace}}

\begin{document}
\thispagestyle{empty}
\begin{singlespace}
\begin{center}
\textbf{ALTO VALOR ADICIONADO NA INDÚSTRIA, NOS SERVIÇOS E NA ADMINISTRAÇÃO PÚBLICA: GUAMARÉ, NATAL E CEARÁ-MIRIM À LUZ DA LITERATURA SOBRE DESENVOLVIMENTO SOCIOECONÔMICO}\footnote{Artigo apresentado à disciplina Estatística Aplicada à Ciência de Dados (POS.1035) do Curso de Especialização em Ciência de Dados do IFRN, Prof. João Maria Filgueira, outubro de 2026.}
\end{center}
\vspace{8pt}
\begin{flushright}
Iuri Guerra de Freitas Pereira\footnote{Pós-graduando em Ciência de Dados pelo Instituto Federal de Educação, Ciência e Tecnologia do Rio Grande do Norte (IFRN). Engenheiro de Computação, Mestre em Tecnologia da Informação. E-mail: iuri.pereira@ifrn.edu.br}
\end{flushright}
\vspace{8pt}
\noindent\textbf{RESUMO}

\noindent O artigo identifica fatores associados ao alto Valor Adicionado Bruto (VAB) da indústria, dos serviços e da administração pública nos municípios do Rio Grande do Norte em 2021, último ano com abertura setorial do VAB municipal, e dialoga com a literatura sobre desenvolvimento socioeconômico. Com dados do IBGE, do IpeaData e do Observatório Sebrae para os 167 municípios, aplicam-se média, desvio padrão, separatrizes (percentis 10 e 90) e a reta $Y = a + bX$; a correlação é apresentada como a inclinação dessa reta quando as variáveis são medidas em desvios padrão. Os casos são Guamaré (indústria), Natal (serviços) e Ceará-Mirim (administração pública). Os 17 municípios abaixo do percentil 10 do PIB per capita têm a administração pública como maior setor e 15 dos 17 acima do percentil 90 têm a indústria, mas o IDHM mediano dos dois grupos é quase o mesmo. O diálogo com Mendes \textit{et al.} (2018) e Fernandes, Tavares e Azevedo (2018) confirma que VAB e PIB caminham juntos e contesta, para o caso potiguar, que a indústria seja o setor mais ligado ao desenvolvimento humano.

\vspace{4pt}
\noindent\textbf{Palavras-chave:} valor adicionado bruto; correlação; desenvolvimento socioeconômico; Rio Grande do Norte.
\end{singlespace}

\vspace{6pt}
\section{Introdução}

Em 2021 o Rio Grande do Norte gerou um Valor Adicionado Bruto de R\$ 71,06 bilhões, dos quais 45,5\% nos serviços, 29,4\% na administração pública, 20,7\% na indústria e 4,4\% na agropecuária (IBGE, 2023). Visto município a município, o estado é outro: a administração pública é o setor de maior VAB em 122 dos 167 municípios (73,1\%), e três cidades (Natal, Mossoró e Parnamirim) concentram 48,8\% do PIB e 63,7\% do VAB dos serviços. A indústria de maior peso está no litoral norte, em municípios pequenos ligados ao petróleo, ao gás, ao sal e à energia eólica (Azevedo; Galindo, 2016).

Este artigo tem como objetivo escolher um município com alto VAB em cada um dos três setores, identificar os fatores dessa realidade e dialogar com a literatura sobre desenvolvimento socioeconômico. O ano de referência é 2021 porque é o último com VAB municipal aberto por setor; as edições de 2022 e 2023 trazem apenas o PIB total, em razão da mudança do ano-base das Contas Nacionais (IBGE, 2025). Os municípios escolhidos são Guamaré, Natal e Ceará-Mirim. A seção 2 apresenta os dados e as ferramentas, incluindo a correlação explicada a partir da relação entre as variáveis; a seção 3 resume os resultados; a seção 4, que é o centro do trabalho, confronta esses resultados com dois artigos; a seção 5 conclui.

\section{Dados e ferramentas}

A base reúne, para os 167 municípios, o PIB e o VAB por setor de 2002 a 2021 (IBGE, tabela 5938 do SIDRA), a população anual, os indicadores sociais e as receitas municipais do IBGE Cidades, o pessoal ocupado do Cadastro Central de Empresas e os estabelecimentos ativos do Observatório Sebrae; os valores a preços de 2010 vêm do IpeaData. O VAB é a diferença entre o valor da produção de uma atividade e o consumo intermediário; o PIB municipal é a soma do VAB das atividades mais os impostos sobre produtos. O IBGE estima o VAB de cada atividade no estado e o reparte entre os municípios por indicadores: quantidade de petróleo produzida, valor das saídas de mercadorias, folha de pagamento das prefeituras (IBGE, 2016). Por isso o grupo ``Indústria'' inclui petróleo, refino, sal e energia elétrica, e o VAB da administração pública dos municípios pequenos reflete a folha da prefeitura.

As ferramentas utilizadas nese artigo para análise são as seguintes: média $\bar{x}$ e o desvio padrão $s$ descrevem o centro e a variação dos dados; a soma dos quadrados dos desvios é mínima quando tomada em relação à média, e $s$ é a raiz dessa soma dividida por $n-1$. As separatrizes dividem os municípios em grupos: o percentil 10 (P10) e o percentil 90 (P90) do PIB per capita isolam os 10\% mais pobres e os 10\% mais ricos. A função afim $Y = a + bX$ relaciona duas variáveis em um gráfico de pontos; seus parâmetros são escolhidos de modo que a soma dos quadrados das distâncias verticais entre cada ponto $(x_i, y_i)$ e o ponto da reta de mesma abscissa, $y_i - (a + b x_i)$, seja a menor possível, o que resulta em
\begin{equation}
b = \frac{\sum (x_i - \bar{x})(y_i - \bar{y})}{\sum (x_i - \bar{x})^2}, \qquad a = \bar{y} - b\,\bar{x}.
\end{equation}

\subsection{Correlação a partir da média, do desvio padrão e da reta}

O coeficiente de correlação de Pearson pode ser entendido sem nenhum conceito além dos já vistos. Primeiro, mede-se cada variável em desvios padrão a partir da sua média: $z_x = (x - \bar{x})/s_x$ e $z_y = (y - \bar{y})/s_y$. Segundo, ajusta-se a reta $Y = a + bX$ a esses valores padronizados. Como as médias de $z_x$ e $z_y$ são zero, a reta passa pela origem ($a = 0$) e sua inclinação é
\begin{equation}
r = \frac{\sum z_{x,i}\, z_{y,i}}{n - 1},
\end{equation}
isto é, a média dos produtos dos valores padronizados. Esse número é o coeficiente de correlação. Ele fica entre $-1$ e $+1$: quando todos os municípios estão do mesmo lado da média nas duas variáveis (ambos acima ou ambos abaixo), os produtos são positivos e $r$ se aproxima de $+1$; quando estão em lados opostos, os produtos são negativos e $r$ se aproxima de $-1$; quando os sinais se misturam, os produtos se cancelam e $r$ fica perto de zero. A ligação com a reta nas unidades originais é imediata: $b = r \cdot s_y / s_x$. Correlação, portanto, é a inclinação da reta $Y = a + bX$ quando $X$ e $Y$ são medidos em desvios padrão.

A Tabela \ref{tab:z} aplica a ideia a seis municípios, com $X$ = participação dos serviços no VAB (\%) e $Y$ = IDHM. As médias são 42,89 e 0,6538 e os desvios padrão 22,39 e 0,0815. Natal está 0,95 desvio padrão acima da média em $X$ e 1,34 acima em $Y$ (produto 1,278); João Dias está 1,05 abaixo em $X$ e 1,52 abaixo em $Y$ (produto 1,589). Todos os produtos são positivos, a soma é 4,067 e $r = 4{,}067/5 = 0{,}81$. Com os 167 municípios o valor cai para $r = 0{,}61$, porque os pontos se espalham mais, mas o sinal se mantém. Dois cuidados: valores muito afastados da média pesam muito no produto, e correlação não é causa.

\begin{table}[H]
\centering
\caption{Correlação como média dos produtos dos valores padronizados: seis municípios ($X$ = participação dos serviços no VAB em 2021, \%; $Y$ = IDHM 2010)}
\label{tab:z}
\footnotesize
\begin{tabular}{@{}lrrrrr@{}}
\toprule
Município & $x_i$ & $y_i$ & $z_x = (x_i - \bar{x})/s_x$ & $z_y = (y_i - \bar{y})/s_y$ & $z_x \cdot z_y$ \\
\midrule
Natal & 64,25 & 0,763 & 0,95 & 1,34 & 1,278 \\
Caicó & 56,57 & 0,710 & 0,61 & 0,69 & 0,421 \\
Pau dos Ferros & 63,51 & 0,678 & 0,92 & 0,30 & 0,273 \\
Ceará-Mirim & 40,38 & 0,616 & $-$0,11 & $-$0,46 & 0,052 \\
Guamaré & 13,17 & 0,626 & $-$1,33 & $-$0,34 & 0,454 \\
João Dias & 19,48 & 0,530 & $-$1,05 & $-$1,52 & 1,589 \\
\midrule
Média & 42,89 & 0,6538 & 0 & 0 & \\
Desvio padrão & 22,39 & 0,0815 & 1 & 1 & Soma = 4,067 \\
\multicolumn{6}{r}{$r = 4{,}067 / (6 - 1) = 0{,}81$} \\
\bottomrule
\end{tabular}
\fonte{Fonte: elaborada pelo autor (2026) com dados do IBGE (2023) e IBGE Cidades.}
\end{table}

\section{Resultados}

A Tabela \ref{tab:p} resume o padrão mais claro dos dados, obtido só com as separatrizes. Os 17 municípios com PIB per capita até o P10 (R\$ 10.039) têm, todos, a administração pública como maior setor, com mediana de 67,1\% do VAB, 6,7\% da população em emprego formal e 93,3\% da receita vinda de transferências. Os 17 municípios a partir do P90 (R\$ 29.137) são, em 15 casos, municípios de indústria, com mediana de 59,4\% do VAB nesse setor. O IDHM mediano, porém, é 0,587 no grupo pobre, 0,608 no grupo do meio e 0,605 no grupo rico: ser muito rico em PIB per capita está associado a ter indústria, não a ter melhor desenvolvimento humano.

\begin{table}[H]
\centering
\caption{Medianas dos grupos definidos pelos percentis 10 e 90 do PIB per capita, 2021}
\label{tab:p}
\footnotesize
\begin{tabular}{@{}p{6.2cm}>{\centering\arraybackslash}p{2.4cm}>{\centering\arraybackslash}p{2.9cm}>{\centering\arraybackslash}p{2.6cm}@{}}
\toprule
Indicador (mediana do grupo) & Até P10 \newline (n = 17) & Entre P10 e P90 \newline (n = 133) & A partir de P90 \newline (n = 17) \\
\midrule
PIB per capita 2021 (R\$) & 9.545 & 13.469 & 47.183 \\
Part. Adm. pública no VAB (\%) & 67,1 & 51,7 & 17,1 \\
Part. Indústria no VAB (\%) & 4,2 & 5,1 & 59,4 \\
Part. Serviços no VAB (\%) & 22,8 & 27,5 & 13,9 \\
Transferências na receita 2025 (\%) & 93,3 & 91,2 & 89,3 \\
\% pop. em emprego formal 2024 & 6,7 & 11,1 & 17,7 \\
IDHM 2010 & 0,587 & 0,608 & 0,605 \\
\bottomrule
\end{tabular}
\fonte{Fonte: elaborada pelo autor (2026) com dados do IBGE (2023), IBGE Cidades e CEMPRE.}
\end{table}

As retas $Y = a + bX$ ajustadas aos 167 municípios resumem as associações em um número cada. A participação da indústria tem correlação $r = +0{,}81$ com o PIB per capita (cada ponto percentual de indústria corresponde a R\$ 968 a mais de PIB per capita), mas $r = -0{,}01$ com o IDHM: a reta é horizontal. A participação dos serviços tem $r = +0{,}61$ com o IDHM ($b = 0{,}00208$, ou 0,021 ponto de IDHM a cada dez pontos percentuais de serviços). A participação da administração pública tem $r = -0{,}57$ com a proporção da população em emprego formal ($b = -0{,}245$) e $r = +0{,}43$ com a dependência de transferências.

Os três municípios escolhidos (Tabela \ref{tab:c}) ilustram os mecanismos. \textbf{Guamaré} (16.261 habitantes) gerou R\$ 1,47 bilhão de VAB industrial, 10,0\% do VAB industrial do estado e 75,0\% do VAB do município; seu PIB per capita de R\$ 125.585 é o terceiro do estado e está 4,6 desvios padrão acima da média, enquanto o IDHM de 0,626 está 0,4 desvio padrão acima, na 43ª posição. Os fatores são a localização do polo de refino e processamento de gás da Bacia Potiguar, operado pela Petrobras até sua venda à 3R Petroleum em 2023 (Silva; Aquino, 2016; 3R Petroleum, 2023), a regra do IBGE que aloca o VAB do petróleo nos municípios produtores e o preço do petróleo, que chegou a tornar o VAB industrial do município negativo entre 2011 e 2014, quando o custo do insumo superou o valor da produção no refino (IBGE, 2015). \textbf{Natal} (896.708 habitantes) concentra 41,3\% do VAB dos serviços do estado; os serviços são 64,2\% do seu VAB, a maior participação entre os 167 municípios, e seu IDHM de 0,763 é o segundo do estado, 4,0 desvios padrão acima da média, embora o PIB per capita (R\$ 26.972, 21ª posição) esteja só 0,3 desvio padrão acima. Os fatores são a função de capital, a densidade, a concentração do comércio, da saúde, da educação e da administração estadual e federal, e a posição de Natal no topo da rede urbana potiguar (Dantas; Clementino, 2014); a preços de 2010, porém, seu PIB per capita atingiu o máximo em 2012 e recuou 27\% até 2021. \textbf{Ceará-Mirim} (74.268 habitantes) tem 46,3\% do VAB na administração pública (R\$ 419 milhões, o maior VAB do setor entre os 122 municípios em que ele domina), indústria de apenas 5,1\%, PIB per capita de R\$ 13.494 (83ª posição) e IDHM de 0,616 (63ª), ambos a menos de 0,3 desvio padrão da média: é o município típico do estado. Os fatores são o declínio da agroindústria canavieira do vale do Ceará-Mirim a partir da década de 1970, sem substituição industrial (Andrade, 2023), e o tamanho da população, que faz da folha pública a maior parcela da economia.

\begin{table}[H]
\centering
\caption{Perfil dos três municípios em 2021 e mediana dos 167 municípios}
\label{tab:c}
\footnotesize
\begin{tabular}{@{}p{4.6cm}>{\centering\arraybackslash}p{2.5cm}>{\centering\arraybackslash}p{2.5cm}>{\centering\arraybackslash}p{2.5cm}>{\centering\arraybackslash}p{2.3cm}@{}}
\toprule
Indicador & Guamaré & Natal & Ceará-Mirim & Mediana RN \\
\midrule
População estimada 2021 & 16.261 & 896.708 & 74.268 & 8.297 \\
PIB per capita 2021 (R\$) e posição & 125.585 (3º) & 26.972 (21º) & 13.494 (83º) & 13.469 \\
Part. no VAB (\%): Indústria / Serviços / Adm. pública & 75,0 / 13,2 / 9,7 & 12,3 / 64,2 / 23,4 & 5,1 / 40,4 / 46,3 & 5,4 / 26,3 / 51,7 \\
IDHM 2010 e posição & 0,626 (43º) & 0,763 (2º) & 0,616 (63º) & 0,606 \\
Afastamento da média (em DP): PIB per capita / IDHM & +4,6 / +0,4 & +0,3 / +4,0 & $-$0,3 / +0,1 & --- \\
\% pop. em emprego formal 2024 & 18,5 & 49,0 & 11,3 & 11,1 \\
Transferências na receita 2025 (\%) & 77,9 & 57,6 & 74,6 & 91,1 \\
\bottomrule
\end{tabular}
\fonte{Fonte: elaborada pelo autor (2026) com dados do IBGE (2023), IBGE Cidades e CEMPRE.}
\end{table}

\section{Diálogo com a literatura}

\subsection{Os dois artigos}

Mendes \textit{et al.} (2018), na Revista de Administração Pública, investigaram a influência da capacidade econômica e da formação de receitas públicas sobre o desenvolvimento humano dos municípios brasileiros entre 2007 e 2013, com um modelo de regressão logística em painel: a variável resposta indica se o município alcança Índice Firjan de Desenvolvimento Municipal (IFDM) de pelo menos 0,600, e as explicativas são o PIB da agropecuária, da indústria e dos serviços, as receitas tributárias próprias e as transferências da União e dos estados. O PIB da indústria teve o maior efeito positivo (razão de chances de 2,55); o PIB dos serviços foi ``a única variável aplicada ao modelo que apresentou efeito negativo''; e as transferências da União foram as variáveis mais relevantes, o que ``pode indicar dependência financeira desses recursos'' (Mendes \textit{et al.}, 2018, p. 930-931). A conclusão é que ``o porte industrial, mensurado pelo PIB indústria, foi o de maior relevância para o desenvolvimento humano'' (Mendes \textit{et al.}, 2018, p. 929).

Fernandes, Tavares e Azevedo (2018), na RACE, estudaram os próprios municípios potiguares: com 166 municípios entre 2010 e 2013, explicaram o PIB municipal pelo VAB da agropecuária, da indústria, dos serviços e da administração pública, além do IFDM e da população. Os quatro VAB tiveram efeito positivo e significativo (coeficientes de 1,39; 1,02; 1,08 e 1,38), o IFDM teve sinal negativo, e os autores concluem que as atividades econômicas influenciam positivamente o PIB e que o PIB mede crescimento, não qualidade de vida.

\subsection{O que os dados de 2021 comprovam}

O resultado de Fernandes, Tavares e Azevedo (2018) é comprovado de forma direta e com as ferramentas mais simples. A participação da indústria tem correlação de $+0{,}81$ com o PIB per capita, os 17 municípios acima do P90 são quase todos industriais e Natal, Mossoró e Parnamirim continuam somando quase metade do PIB estadual. Há uma ressalva: parte dessa relação é contábil, porque o PIB é, por definição, a soma dos VAB mais os impostos; o que os dados acrescentam é a magnitude, com a indústria elevando o PIB per capita muito mais do que os outros setores. O sinal negativo que os autores encontraram para o IFDM também é coerente com a Tabela \ref{tab:p}: os municípios mais ricos em PIB per capita não têm IDHM maior que os demais. A aparente contradição entre VAB alto e desenvolvimento humano baixo, que os autores registram sem explicar, ganha aqui um mecanismo: nos municípios de enclave, o VAB é contabilizado no território, mas o lucro, os salários mais altos e os impostos federais são apropriados fora dele. Silva e Aquino (2016) já haviam mostrado, com dados de 2000 a 2011, que Guamaré tinha o maior PIB per capita do estado e royalties equivalentes a 3,49 vezes a cota do FPM, mas IDHM abaixo da média estadual, e chamaram o fenômeno de ``paradoxo do subdesenvolvimento com abundância de divisas''. Os dados de 2021 mostram que o paradoxo persiste dez anos depois e que se estendeu aos municípios eólicos: Parazinho e São Bento do Norte, com 88,7\% do VAB na indústria e PIB per capita de R\$ 114 mil e R\$ 204 mil, têm IDHM de 0,549 e 0,555, entre os seis piores do estado.

O resultado de Mendes \textit{et al.} (2018) sobre as transferências também é comprovado, desde que se distinga volume de dependência. As transferências sustentam as prefeituras de quase todos os municípios potiguares (mediana de 91\% da receita), e por isso aparecem no modelo nacional como o fator mais associado a um IFDM razoável. Mas a dependência mais alta está no grupo abaixo do P10 (93,3\%), onde o emprego formal é menor (6,7\%) e o IDHM é o pior (0,587). É o retrato que Gomes e Mac Dowell (2000) anteciparam ao mostrar que a criação de municípios e o aumento das transferências após 1988 beneficiaram os pequenos municípios sem estimular sua atividade produtiva, e que Araújo e Lima (2009) chamaram de ``economia sem produção'': renda formada por salários públicos, aposentadorias e transferências, não por produção. Ceará-Mirim, com 46\% do VAB na administração pública, e os 122 municípios do mesmo perfil são a expressão potiguar desse diagnóstico.

\subsection{O que os dados de 2021 contestam}

A conclusão de Mendes \textit{et al.} (2018) de que a indústria é o setor mais relevante para o desenvolvimento humano não se sustenta no Rio Grande do Norte, e a razão aparece nas separatrizes. No Brasil como um todo, PIB industrial alto costuma significar indústria de transformação em cidades grandes, com empregos, fornecedores locais e serviços associados, que é o que a literatura chama de encadeamentos. No Rio Grande do Norte, a participação industrial extrema ocorre em municípios pequenos com uma única planta de petróleo, sal ou energia eólica: o grupo acima do P90 tem mediana de 10.441 habitantes, 59,4\% do VAB na indústria e IDHM mediano de 0,605, igual ao do meio da distribuição. Parte da diferença é de medida: os autores usam o PIB setorial em nível, que cresce com o tamanho do município e por isso carrega o efeito das cidades grandes, enquanto este trabalho usa a participação de cada setor, que mede especialização e não depende do tamanho. A outra parte é de conteúdo: o efeito da indústria sobre o desenvolvimento humano depende do tipo de indústria, e a indústria potiguar de maior VAB é extrativa e energética, sem os encadeamentos que o resultado nacional pressupõe.

O sinal negativo dos serviços encontrado por Mendes \textit{et al.} (2018) também é contestado. Os autores o atribuem à presença de serviços de baixo valor agregado em todos os municípios, inclusive nos pequenos. Nos dados potiguares, a participação dos serviços é a variável mais fortemente associada ao IDHM ($r = +0{,}61$) e à menor pobreza de renda, e Natal, com 64\% do VAB em serviços, tem o segundo melhor IDHM do estado. A explicação provável está na classificação: no Rio Grande do Norte, o terciário de baixo valor agregado dos municípios pequenos aparece como administração pública, que o IBGE separa dos serviços, e não como serviços mercantis. Quando a administração pública é tratada à parte, o sinal dos serviços se inverte; quando é somada aos serviços, como em muitos estudos nacionais, o peso dos municípios pequenos puxa o coeficiente para baixo.

Em resumo, o que acontece em um município pode ser referenciado em outro de mesmo tipo (Guamaré repete o padrão dos municípios petrolíferos de Silva e Aquino, 2016; Ceará-Mirim, o dos municípios dependentes de transferências de Araújo e Lima, 2009; Natal, o das capitais na rede urbana de Dantas e Clementino, 2014), mas não entre tipos diferentes, distinção que um modelo nacional com todos os municípios não capta.

\section{Considerações finais}

Com média, desvio padrão, separatrizes e a reta $Y = a + bX$, foi possível identificar os fatores do alto VAB de Guamaré (polo de petróleo e gás, regra de alocação do VAB e preço do petróleo), de Natal (função de capital, densidade e concentração de comércio, saúde, educação e administração) e de Ceará-Mirim (declínio da cana sem substituição industrial e folha pública como base da economia), e confrontá-los com a literatura. O coeficiente de correlação, apresentado como a inclinação da reta entre variáveis medidas em desvios padrão, resumiu as relações: indústria e PIB per capita sobem juntos ($+0{,}81$), indústria e IDHM não ($-0{,}01$), serviços e IDHM sobem juntos ($+0{,}61$), administração pública e emprego formal caminham em sentidos opostos ($-0{,}57$). O diálogo com Fernandes, Tavares e Azevedo (2018) e Mendes \textit{et al.} (2018) confirmou que VAB e PIB caminham juntos e que as transferências sustentam os municípios pequenos, e contestou, para o caso potiguar, que a indústria seja o setor mais ligado ao desenvolvimento humano e que os serviços o prejudiquem. A lição metodológica é que a medida escolhida (nível ou participação) e a classificação dos setores (serviços com ou sem administração pública) mudam o sinal das conclusões, o que recomenda cautela ao transferir resultados nacionais para um estado cuja indústria é de enclave. Trabalhos futuros devem repetir a análise quando o IBGE retomar a abertura setorial do VAB municipal, prevista para 2027, e quando o IDHM do Censo de 2022 for divulgado.

\section*{REFERÊNCIAS}
\begin{referencias}

3R PETROLEUM ÓLEO E GÁS S.A. \textbf{Fato relevante}: conclusão da aquisição do Polo Potiguar. Rio de Janeiro: 3R Petroleum, 7 jun. 2023. Disponível em: \url{https://api.mziq.com/mzfilemanager/v2/d/55b913af-cd4c-48d5-bc19-48c63916b8a5/56bf8b8b-7ebd-c7dd-e0a9-f31d07f5c7da}. Acesso em: 10 out. 2026.

ANDRADE, Manuel Correia de. As usinas de açúcar e as destilarias de álcool do Rio Grande do Norte. \textbf{Revista GeoInterações}, Assú, v. 7, e11, p. 175-185, 2023. Disponível em: \url{https://periodicos.apps.uern.br/index.php/RGI/article/view/5500}. Acesso em: 10 out. 2026.

ARAÚJO, Leonardo Alves de; LIMA, João Policarpo Rodrigues. Transferências de renda e empregos públicos na economia sem produção do semiárido nordestino. \textbf{Planejamento e Políticas Públicas}, Brasília, n. 33, p. 45-77, jul./dez. 2009. Disponível em: \url{https://repositorio.ipea.gov.br/bitstream/11058/3993/2/PPP_n33_Transferencias.pdf}. Acesso em: 10 out. 2026.

AZEVEDO, Francisco Fransualdo de; GALINDO, Leonardo da Silva. A espacialidade da indústria no Rio Grande do Norte (Brasil) no contexto da reestruturação produtiva. \textbf{Boletim Gaúcho de Geografia}, Porto Alegre, v. 43, n. 1, p. 147-168, 2016. Disponível em: \url{https://seer.ufrgs.br/index.php/bgg/article/view/57933}. Acesso em: 10 out. 2026.

DANTAS, Joseney Rodrigues de Queiroz; CLEMENTINO, Maria do Livramento Miranda. Desenvolvimento e urbanização: a rede urbana potiguar sob a ótica dos ``REGICs''. \textbf{Geotemas}, Pau dos Ferros, v. 4, n. 1, p. 71-82, jan./jun. 2014. Disponível em: \url{https://periodicos.apps.uern.br/index.php/GEOTemas/article/view/606}. Acesso em: 10 out. 2026.

FERNANDES, Ruhama Bezerra; TAVARES, Adilson de Lima; AZEVEDO, Yuri Gomes Paiva. Relação entre o valor adicionado das atividades econômicas e o Produto Interno Bruto do Rio Grande do Norte. \textbf{RACE: Revista de Administração, Contabilidade e Economia}, Joaçaba, v. 17, n. 2, p. 757-782, maio/ago. 2018. DOI: 10.18593/race.v17i2.17238.

GOMES, Gustavo Maia; MAC DOWELL, Maria Cristina. \textbf{Descentralização política, federalismo fiscal e criação de municípios}: o que é mau para o econômico nem sempre é bom para o social. Brasília: IPEA, 2000. (Texto para Discussão, n. 706). Disponível em: \url{http://repositorio.ipea.gov.br/handle/11058/2339}. Acesso em: 10 out. 2026.

IBGE. \textbf{Produto interno bruto dos municípios 2013}: apresentação da divulgação. Rio de Janeiro: IBGE, Coordenação de Contas Nacionais, 18 dez. 2015. Disponível em: \url{https://anda.ibge.gov.br/media/com_mediaibge/arquivos/d7449516f181496dbff2a68cfdd9c349.pdf}. Acesso em: 10 out. 2026.

IBGE. \textbf{Produto interno bruto dos municípios}: ano de referência 2010. 3. ed. Rio de Janeiro: IBGE, 2016. (Relatórios Metodológicos, v. 29). Disponível em: \url{https://biblioteca.ibge.gov.br/visualizacao/livros/liv97483.pdf}. Acesso em: 10 out. 2026.

IBGE. \textbf{Produto interno bruto dos municípios 2021}: notas técnicas. Rio de Janeiro: IBGE, 2023. (Contas Nacionais, n. 96). Disponível em: \url{https://biblioteca.ibge.gov.br/visualizacao/livros/liv102057_notas_tecnicas.pdf}. Acesso em: 10 out. 2026.

IBGE. \textbf{Produto interno bruto dos municípios}: Brasil: 2022-2023: notas técnicas. Rio de Janeiro: IBGE, 2025. (Contas Nacionais, n. 104). Disponível em: \url{https://biblioteca.ibge.gov.br/visualizacao/livros/liv102250_notas_tecnicas.pdf}. Acesso em: 10 out. 2026.

MENDES, Wesley de Almeida; FERREIRA, Marco Aurélio Marques; ABRANTES, Luiz Antônio; FARIA, Evandro Rodrigues de. A influência da capacidade econômica e da formação de receitas públicas no desenvolvimento humano. \textbf{Revista de Administração Pública}, Rio de Janeiro, v. 52, n. 5, p. 918-934, set./out. 2018. DOI: 10.1590/0034-761220170004.


SILVA, Maria Russilânya Costa Santos da; AQUINO, Joacir Rufino de. Subdesenvolvimento humano com abundância de riqueza: o caso dos ``municípios petrorrentistas'' do Rio Grande do Norte. \textbf{Revista GeoNordeste}, São Cristóvão, ano 27, n. 1, p. 71-91, jan./jul. 2016. Disponível em: \url{https://periodicos.ufs.br/geonordeste/article/view/4572}. Acesso em: 10 out. 2026.

\end{referencias}
\end{document}
